\documentclass[12pt]{article}

\usepackage[paperwidth=612pt,paperheight=792pt,top=72pt,bottom=72pt,left=30mm,right=20mm]{geometry}
\usepackage{times}
\usepackage{setspace}
\usepackage{fancyhdr}
\usepackage{ragged2e}
\usepackage{booktabs}
\setstretch{1.5}

\pagestyle{fancy}
\fancyhf{}
\fancyhead[R]{AyurSpace}
\fancyfoot[L]{Pillai HOC College of Engineering \& Technology (Autonomous), B. E. Computer Engineering}
\fancyfoot[R]{\thepage}
\renewcommand{\headrulewidth}{1pt}
\renewcommand{\footrulewidth}{1pt}

\begin{document}

\vspace*{\fill}
\begin{center}
\textbf{\Huge{Chapter 2}}\\
\vspace{5mm}
\textbf{\Huge{Literature Review}}
\end{center}
\vspace*{\fill}
\newpage

\section*{2.1 \quad Related Work}
\addcontentsline{toc}{section}{2.1 Related Work}
\justify

\noindent\textbf{Computer Vision in Botany:}\\
The area of automated plant identification has grown significantly with the advent of Convolutional
Neural Networks (CNNs). Initial endeavors employed leaf-shape descriptors and edge detection
algorithms [Waldchen et al., 2018]. Modern methods use Deep Learning architectures such as
ResNet-50 and MobileNetV3. For instance, \textbf{Pl@ntNet}, a citizen-science project, uses a huge
shared dataset that does a great job with European plants but has trouble with Indian medicinal
herbs. Similarly, \textbf{LeafSnap} focuses on leaf curvature to determine tree species
characteristics. The main limitation with these systems is that they work in a ``Botanical Silo.''
They provide a Latin binomial---the end of the interaction for them---whereas for Ayurveda,
identification is just the beginning of the consultation.

\vspace{5mm}

\noindent\textbf{LLMs in Healthcare:}\\
Generative AI has shown promise in combining medical knowledge. For instance, Singhal et al.
[2023] showed that Med-PaLM could pass the USMLE. But applying general-purpose LLMs to specific
areas of traditional knowledge bases frequently leads to ``alignment drift,'' in which the model
prioritizes Western medical interpretations ahead of traditional logic due to bias in the training
data. There is no dedicated ``Ayur-PaLM,'' leading to a gap where prompt engineering methods are
necessary to limit general LLMs (like Google Gemini) to act as reliable experts in Ayurveda.

\vspace{5mm}

\noindent\textbf{Digital Approaches to Ayurveda:}\\
Previous research in ``Computational Ayurveda'' has concentrated primarily on structural diagnostics
and static databases. Studies include \textbf{Pulse Diagnosis (\textit{Nadi Pariksha})}, using
piezoelectric sensors to turn pulse waveforms into numbers, and \textbf{Tongue Diagnosis
(\textit{Jivha Pariksha})}, using image processing to find color and coating. Further, existing
``Ayur Apps'' provide simple static databases---CRUD apps serving as basic electronic dictionaries
without active AI context or personalized adaptability.

\vspace{5mm}

\begin{table}[htbp]
\caption{AyurSpace Compared to Other Solutions}
\begin{center}
\begin{tabular}{|c|c|c|c|c|}
\hline
\textbf{Feature} & \textbf{PlantNet} & \textbf{Google Lens} & \textbf{Ayur Apps} & \textbf{AyurSpace} \\
\hline
\textbf{Visual ID} & High & High & None & \textbf{High} \\
\hline
\textbf{Context} & None & Limited & Static & \textbf{Dynamic} \\
\hline
\textbf{Dosha Logic} & None & None & Basic & \textbf{Adaptive} \\
\hline
\textbf{Interaction} & Passive & Passive & Passive & \textbf{Chat} \\
\hline
\textbf{Safety} & N/A & Low & Medium & \textbf{High} \\
\hline
\end{tabular}
\label{tab1}
\end{center}
\end{table}

AyurSpace bridges the gap between high-tech vision and high-touch traditional wisdom.

\end{document}
